\subsection{Condicionamento Numérico}

\pgfplotstableread[col sep=comma]{../../results/ill_conditioned_gram.csv}\illconditioned

\pgfplotstablegetelem{0}{num_points}\of\illconditioned
\edef\illHighPoints{\pgfplotsretval}
\pgfplotstablegetelem{0}{degree}\of\illconditioned
\edef\illHighDegree{\pgfplotsretval}
\pgfplotstablegetelem{0}{cond}\of\illconditioned
\edef\illHighCondRaw{\pgfplotsretval}
\pgfplotstablegetelem{1}{ridge_effective}\of\illconditioned
\edef\illHighRidge{\pgfplotsretval}

\pgfplotstablegetelem{2}{num_points}\of\illconditioned
\edef\illDupPoints{\pgfplotsretval}
\pgfplotstablegetelem{2}{degree}\of\illconditioned
\edef\illDupDegree{\pgfplotsretval}
\pgfplotstablegetelem{2}{cond}\of\illconditioned
\edef\illDupCondRaw{\pgfplotsretval}

\edef\illHighCond{\fpeval{\illHighCondRaw}}

Experimentos controlados confirmaram que certas configurações tornam o bloco de
Gram extremamente mal condicionado.
Um script dedicado\footnote{Disponível no repositório como
\texttt{experiments/run\_ill\_conditioned\_regression.py}, o qual gera a
tabela \texttt{results/ill\_conditioned\_gram.csv}.}
reproduz dois cenários representativos: (i)
\illHighPoints\,amostras equiespaçadas em $[-1,1]$ ajustadas com um polinômio
de grau \illHighDegree, cuja matriz apresenta
$\operatorname{cond}(G) \approx
\num[scientific-notation=true,round-mode=figures,round-precision=3]{\illHighCond}$;
e (ii) \illDupPoints\,amostras coincidentes sob ajuste de grau
\illDupDegree, para as quais $\operatorname{cond}(G)$ é reportado como
\texttt{\illDupCondRaw}, evidenciando a degenerescência da matriz de
Vandermonde.

Denotemos por $V \in \mathbb{R}^{N \times p}$ a matriz de
Vandermonde cujas linhas são os vetores de monômios avaliados em cada ponto
$p_j$ e cujas colunas correspondem às $p=\binom{n+r}{r}$ funções base.
Ao acumular $V^T V$ ao longo das amostras, obtemos precisamente o bloco de
Gram $G$, enquanto a acumulação $V^T \vec{f}$ gera o vetor $\vec{b}$.
Quando as colunas de $V$ ficam quase linearmente dependentes - como ocorre em
graus altos, em domínios muito alongados ou em variáveis fortemente
correlacionadas - essa acumulação herda a perda de significância e culmina em
um $G$ mal condicionado.

Para mitigar tais problemas, o solver \texttt{regression\_optimized} foi projetado para identificar automaticamente situações de mal condicionamento e, nesses casos,
ele substitui a resolução em blocos por uma chamada a
\texttt{numpy.linalg.lstsq}, enquanto \texttt{regression\_largeP} mantém a
decomposição, porém injeta automaticamente uma regularização de magnitude
\num[scientific-notation=true]{\illHighRidge} no sistema, conforme registrado
para ambos os cenários analisados.
